\section*{Abstract}

Reducing noise and vibration in lightweight structures is a common engineering challenge. Plates and panels in cars, aircraft and machines often need extra mass or damping to reduce the noise they radiate, and a lightweight way to provide it is to attach a few small tuned mass resonators to the plate. Each resonator is a small spring, damper and mass combination placed at a position on the plate, and together these resonators determine how much sound power the plate radiates. This thesis describes that radiated sound with the equivalent radiated power (ERP), evaluated between 10 and 160 Hz.

Finding a good resonator layout by hand is hard, because the number of resonators, their values and their positions all interact. An exact, closed-form solution exists only for a few simple cases; most realistic scenarios need a full multiphysics simulation, which is accurate but far too slow to try out thousands of candidate layouts, and it only answers one direction of the problem: from a design to its spectrum. What an engineer actually wants is the other direction too: given a target noise behaviour, propose a resonator design that achieves it. This is called generative, or inverse, design, and it is the main goal of this thesis.

To reach that goal, a dataset of resonator designs and their ERP spectra is generated with a validated plate solver. Each resonator is described by five variables, mass, tuning frequency, position and a damping ratio, sampled with Latin hypercube sampling so the dataset covers the design space evenly. Ten neural network architectures are trained as fast forward surrogates on a two-resonator version of this dataset, learning to predict the ERP spectrum of a design in a fraction of a second instead of the time a full simulation takes; the best of them reaches about 1.1\,dB root-mean-square error on unseen designs. The same ten architectures are also compared on a three-resonator dataset, where the error increases by a factor of 1.1 to 2.0 depending on the architecture (for the best one from about 1.1\,dB to 2.2\,dB), showing how much harder the problem becomes as more resonators are added. Building on the forward step, invertible neural operators are developed that use one set of weights for both directions, so the same model can predict a spectrum from a design and, in the other direction, propose a design for a target spectrum. These are evaluated on an earlier, larger two-resonator dataset and already recover several design variables reasonably well, though recovering the exact resonator positions remains difficult, since several different layouts can produce nearly the same spectrum. A separate experiment also checks how well the forward models extrapolate to frequencies they never saw during training, which turns out to be considerably harder than interpolating within the training range.

\newpage
